% ── 1. Introduction ───────────────────────────────────────────────────────────
\section{Introduction}

Large language models (LLMs) have transformed what is possible in natural
language processing, yet they remain unreliable when a task demands precise,
verifiable knowledge~\cite{lewis2020retrieval}. The core problem is that
LLMs store knowledge implicitly in their weights---there is no mechanism to
check what they know, no way to update it without retraining, and no
guarantee that what they generate is grounded in fact.
Retrieval-Augmented Generation (RAG) addresses this by decoupling knowledge
from the model~\cite{gao2023retrieval}: at query time, a retrieval system pulls
relevant documents from an external corpus and passes them as context, so the
LLM reasons over retrieved evidence rather than recalled parameters.
This is especially important in knowledge-intensive domains---legal analysis,
scientific literature review, clinical decision support---where a wrong or
unsupported answer carries real consequences.

Early RAG systems retrieve flat text chunks by vector similarity.
For many questions this is adequate, but it breaks down when the answer
requires connecting facts that are distributed across the corpus rather than
co-located in a single passage.
Graph-based RAG systems~\cite{edge2024graphrag,guo2024lightrag,
gutierrez2024hipporag,chen2025pathrag} address this by building a knowledge
graph (KG) from the corpus and retrieving coherent subgraphs.
The difference is concrete: the question
\textit{``What connects Varroa Destructor to Polyface Farm?''}
has no answer in any single passage, but the knowledge graph contains a
path of three edges through two intermediate nodes (Varroa Destructor
$\to$ Beekeeping Industry $\to$ Farmers $\to$ Polyface) that a graph
traversal can find and a vector search cannot.

Despite this progress, every graph-RAG system we are aware of shares a
critical limitation: it applies one fixed retrieval strategy to every query.
LightRAG~\cite{guo2024lightrag} always expands a 1-hop neighborhood.
HippoRAG~\cite{gutierrez2024hipporag} always runs Personalized PageRank.
PathRAG~\cite{chen2025pathrag} always does flow-based path retrieval.
None of them read the question carefully enough to know which strategy is
appropriate.

This matters because the question itself already tells you what kind of
answer it needs~\cite{li2002learning}.
\textit{``What is Varroa Mite?''} asks for a self-contained definition: one
entity, its description, its immediate connections.
\textit{``What connects Rational Grazing to Polyface Farm?''} asks for a
mechanism: an intermediate entity and the chain of relations linking it to
both anchors~\cite{yang2018hotpotqa}.
\textit{``How do Varroa Mite, Colony Collapse Disorder, and Neonicotinoids
interact?''} asks for synthesis across three anchors simultaneously~\cite{ho2020constructing}.
\textit{``What are the major sustainability themes in modern farming?''} asks
for a thematic sweep with no specific entity anchor at all.
These four shapes are not arbitrary categories---they correspond directly to
four distinct graph traversal algorithms.
A system that ignores these linguistic cues and applies one algorithm to all
four risks retrieving a mismatched subgraph structure for many queries.

The failure is not hypothetical.
PathRAG, applied to
\textit{``What connects Rational Grazing to Polyface?''}, produces a
fluent paragraph about a ``shared commitment to sustainable practices''
without ever naming the intermediate entity---because its flow-based pruning
assigns high reliability to the most-traversed edges (broad category links),
not to the specific low-frequency bridge edge the question is asking about.
This is not a retrieval error; the edge exists in the graph.
It is a routing error: the wrong algorithm was applied to a question that
called for a different one.

We present \textbf{GRASP-RAG} (Graph Retrieval with Adaptive Structured
Paths).
The core idea is simple: read the query's linguistic shape before deciding
how to traverse the graph.
An LLM classifier maps each query to one of four retrieval \emph{motifs}:
an ego-network for definitional questions, a bidirectional shortest-path
traversal for relational chain questions, a Steiner tree approximation for
multi-entity synthesis questions, and LightRAG hybrid mode for broad
thematic questions.
Anchor nodes are extracted from VDB candidates rather than by fuzzy string
matching, reducing entity-linking errors from name paraphrasing; when the
correct entity falls outside the top-20 candidates, anchor recall suffers
instead.
Each retrieved subgraph is then linearized into a topology-aware context
that presents the evidence in the structure the question demands---a labeled
step-by-step chain for bridge queries, an entity profile for single-entity
queries, a comparison scaffold for multi-entity queries.
When the classifier is uncertain or the retrieved subgraph is too sparse, a
confidence score triggers a graceful fallback to LightRAG hybrid rather than
a hard failure.

Prior work has approached adjacent problems without solving this one.
PathRAG introduces flow-based path retrieval but applies the same path
format to every query type.
HippoRAG's Personalized PageRank diffuses too broadly for entity-centric
questions.
Layered Query Retrieval~\cite{huang2024layered} classifies queries by
complexity (simple vs.\ complex), but this binary split cannot distinguish
an ego-network query from a bridge traversal query, and it operates on flat
text without a knowledge graph.
We discuss these and related approaches in detail in
Section~\ref{sec:related}.

GRASP-RAG contributes two things:
\begin{enumerate}
  \item \textbf{Query-adaptive motif selection}: a lightweight LLM classifier
    maps each query to the matching graph traversal algorithm, with
    candidate-constrained anchor extraction that reduces entity-linking
    errors from name paraphrasing, and a confidence-based fallback that limits
    failure on uncertain queries.
  \item \textbf{Topology-aware context structuring}: each motif serializes its
    retrieved subgraph into a structured context that mirrors the answer's
    reasoning shape, so the LLM receives evidence in the form the question
    requires---not a flat list of passages regardless of query type.
\end{enumerate}

Evaluated on four domains of UltraDomain~\cite{ding2024ultradomain}
(Agriculture, CS, Legal, Health) using reference-anchored multi-judge
evaluation, GRASP-RAG achieves consistent win rates against all three
baselines, with an average of 73.4\% against LightRAG and 79.7\% against
PathRAG, the strongest baseline.
The system's advantage holds across query types even when the classifier
makes errors---a robustness finding we examine alongside the main results in
Section~\ref{sec:results}.
